\documentclass{beamer}

%================ ENCODING =================
\usepackage[utf8]{inputenc}
\usepackage[T1]{fontenc}
\usepackage[vietnamese]{babel}

%================ PACKAGES =================
\usepackage{tikz}
\usetikzlibrary{arrows.meta,positioning}
\usepackage[most]{tcolorbox}
\usepackage{amsmath}
\usepackage{xcolor}
\usepackage{fontawesome5}

%================ COLORS (PASTEL LIGHT) =================
\definecolor{softblue}{RGB}{173,216,255}
\definecolor{softpurple}{RGB}{210,180,255}
\definecolor{bg1}{RGB}{245,245,255}
\definecolor{bg2}{RGB}{230,220,255}

%================ BACKGROUND =================
\usebackgroundtemplate{
\begin{tikzpicture}[remember picture,overlay]
\shade[left color=bg1,right color=bg2]
(current page.south west) rectangle (current page.north east);

\fill[softblue,opacity=0.2] (current page.north east) circle (2cm);
\fill[softpurple,opacity=0.2] (current page.south west) circle (2cm);
\end{tikzpicture}
}

%================ STYLE =================
\setbeamercolor{normal text}{fg=black}
\setbeamertemplate{navigation symbols}{}

%================ FOOTLINE =================
\setbeamercolor{footline}{fg=black,bg=white}
\setbeamertemplate{footline}{
\leavevmode
\hbox{
\begin{beamercolorbox}[wd=0.5\paperwidth,ht=2.5ex,dp=1ex,left]{footline}
\hspace{0.3cm} Hoàng Anh Tú
\end{beamercolorbox}
\begin{beamercolorbox}[wd=0.5\paperwidth,ht=2.5ex,dp=1ex,right]{footline}
\insertframenumber{} / \inserttotalframenumber \hspace{0.3cm}
\end{beamercolorbox}
}
}

%================ BOX =================
\tcbset{
colback=white,
colframe=softblue!70!black,
arc=10pt,
boxrule=0.8pt
}

\begin{document}
%================ SLIDE: KHÁI NIỆM ĐỒ THỊ =================
\begin{frame}{\faProjectDiagram\ Khái niệm đồ thị}
\begin{tcolorbox}
\begin{itemize}
\item Đồ thị là cấu trúc toán học: $G = (V, E)$
\item $V$ là tập các đỉnh (vertices)
\item $E$ là tập các cạnh (edges)
\item Mỗi cạnh nối hai đỉnh biểu diễn quan hệ
\item Có thể là đồ thị vô hướng hoặc có hướng
\end{itemize}
\end{tcolorbox}
\end{frame}

%================ SLIDE: THÀNH PHẦN =================
\begin{frame}{\faCubes\ Các thành phần cơ bản}
\begin{tcolorbox}
\begin{itemize}
\item Đỉnh (Vertex): điểm trong đồ thị
\item Cạnh (Edge): kết nối giữa hai đỉnh
\item Bậc đỉnh: số cạnh liên kết với đỉnh
\item Đường đi (Path): dãy các đỉnh liên tiếp
\item Chu trình (Cycle): đường đi khép kín
\end{itemize}
\end{tcolorbox}
\end{frame}

%================ SLIDE: VÍ DỤ =================
\begin{frame}{\faLightbulb\ Ví dụ minh họa}
\begin{tcolorbox}
\begin{itemize}
\item Đồ thị gồm 4 đỉnh: 1, 2, 3, 4
\item Các cạnh: (1,2), (1,3), (2,4)
\item Biểu diễn quan hệ giữa các đối tượng
\item Có thể là mạng bạn bè
\item Hoặc sơ đồ đường đi
\end{itemize}
\end{tcolorbox}

\vspace{0.3cm}

\begin{center}
\begin{tikzpicture}[scale=1]
\node[circle,draw,fill=softblue!60] (1) at (0,0) {1};
\node[circle,draw,fill=softblue!60] (2) at (2,0) {2};
\node[circle,draw,fill=softpurple!60] (3) at (1,1.5) {3};
\node[circle,draw,fill=softpurple!60] (4) at (1,-1.5) {4};

\draw (1)--(2);
\draw (1)--(3);
\draw (2)--(4);
\end{tikzpicture}
\end{center}
\end{frame}

%================ SLIDE: ỨNG DỤNG =================
\begin{frame}{\faNetworkWired\ Ứng dụng của đồ thị}
\begin{tcolorbox}
\begin{itemize}
\item Mạng xã hội (Facebook, Instagram)
\item Bản đồ và hệ thống GPS
\item Mạng máy tính, Internet
\item Trí tuệ nhân tạo và tìm kiếm
\item Bài toán đường đi ngắn nhất
\end{itemize}
\end{tcolorbox}
\end{frame}

%================ SLIDE: PHÂN LOẠI ĐỒ THỊ (CÓ HÌNH) =================
\begin{frame}{\faSitemap\ Phân loại đồ thị}
\begin{tcolorbox}
\begin{itemize}
\item Đồ thị vô hướng (Undirected)
\item Đồ thị có hướng (Directed)
\item Đồ thị có trọng số (Weighted)
\item Đồ thị đầy đủ (Complete)
\item Đồ thị liên thông (Connected)
\end{itemize}
\end{tcolorbox}

\vspace{0.3cm}

\begin{columns}

%----------------- VÔ HƯỚNG -----------------
\column{0.33\textwidth}
\centering
\textbf{Vô hướng}
\begin{tikzpicture}[scale=0.7]
\node[circle,draw,fill=softblue!60] (A) at (0,0) {A};
\node[circle,draw,fill=softblue!60] (B) at (1.5,0) {B};
\node[circle,draw,fill=softblue!60] (C) at (0.75,1.3) {C};
\draw (A)--(B);
\draw (B)--(C);
\draw (C)--(A);
\end{tikzpicture}

%----------------- CÓ HƯỚNG -----------------
\column{0.33\textwidth}
\centering
\textbf{Có hướng}
\begin{tikzpicture}[scale=0.7,>=stealth]
\node[circle,draw,fill=softpurple!60] (A) at (0,0) {A};
\node[circle,draw,fill=softpurple!60] (B) at (1.5,0) {B};
\node[circle,draw,fill=softpurple!60] (C) at (0.75,1.3) {C};
\draw[->] (A)--(B);
\draw[->] (B)--(C);
\draw[->] (C)--(A);
\end{tikzpicture}

%----------------- TRỌNG SỐ -----------------
\column{0.33\textwidth}
\centering
\textbf{Trọng số}
\begin{tikzpicture}[scale=0.7]
\node[circle,draw,fill=softblue!60] (A) at (0,0) {A};
\node[circle,draw,fill=softblue!60] (B) at (1.5,0) {B};
\node[circle,draw,fill=softblue!60] (C) at (0.75,1.3) {C};
\draw (A)-- node[above]{2} (B);
\draw (B)-- node[right]{3} (C);
\draw (C)-- node[left]{1} (A);
\end{tikzpicture}

\end{columns}

\vspace{0.4cm}

\begin{center}
\textit{Minh họa trực quan các loại đồ thị cơ bản}
\end{center}

\end{frame}
%================ SLIDE: KÝ HIỆU TOÁN HỌC =================
\begin{frame}{\faSquareRootAlt\ Ký hiệu và tính chất}
\begin{tcolorbox}
\begin{itemize}
\item $|V|$ là số lượng đỉnh
\item $|E|$ là số lượng cạnh
\item $\deg(v)$ là bậc của đỉnh $v$
\item $\sum \deg(v) = 2|E|$ (đồ thị vô hướng)
\item Áp dụng trong phân tích cấu trúc đồ thị
\end{itemize}
\end{tcolorbox}
\end{frame}

%================ SLIDE: ĐƯỜNG ĐI =================
\begin{frame}{\faRoute\ Các loại đường đi}
\begin{tcolorbox}
\begin{itemize}
\item Path: dãy đỉnh liên tiếp
\item Simple Path: không lặp đỉnh
\item Cycle: đường đi khép kín
\item Euler Path: đi qua mỗi cạnh đúng 1 lần
\item Hamilton Path: đi qua mỗi đỉnh đúng 1 lần
\end{itemize}
\end{tcolorbox}
\end{frame}

%================ SLIDE: TÍNH CHẤT =================
\begin{frame}{\faCheckCircle\ Tính chất đồ thị}
\begin{tcolorbox}
\begin{itemize}
\item Đồ thị liên thông hay không
\item Có tồn tại chu trình
\item Số thành phần liên thông
\item Bậc của các đỉnh
\item Ảnh hưởng đến thuật toán duyệt
\end{itemize}
\end{tcolorbox}
\end{frame}

%================ SLIDE: MÔ HÌNH HÓA =================
\begin{frame}{\faNetworkWired\ Mô hình hóa thực tế}
\begin{tcolorbox}
\begin{itemize}
\item Người dùng → đỉnh
\item Quan hệ → cạnh
\item Mạng xã hội (Facebook)
\item Bản đồ giao thông (GPS)
\item Mạng máy tính và Internet
\end{itemize}
\end{tcolorbox}
\end{frame}

%================ SLIDE: LIÊN HỆ THUẬT TOÁN =================
\begin{frame}{\faCogs\ Liên hệ thuật toán}
\begin{tcolorbox}
\begin{itemize}
\item DFS: duyệt theo chiều sâu
\item BFS: duyệt theo chiều rộng
\item BFS tìm đường ngắn nhất (không trọng số)
\item DFS dùng kiểm tra chu trình
\item Cơ sở cho nhiều thuật toán nâng cao
\end{itemize}
\end{tcolorbox}
\end{frame}

\end{document}
